\clearpage
\VerdierPhysicalPageAnchor{0186}%
\begingroup\small\setlength{\parskip}{0pt}\setlength{\emergencystretch}{1.5em}
\setlength{\abovedisplayskip}{4pt plus 2pt minus 2pt}%
\setlength{\belowdisplayskip}{4pt plus 2pt minus 2pt}%

\VerdierSemanticAnchor{ch3d-2}
\section*{2. 分解。}

\noindent\textit{在本节中，$\mathcal A$ 表示一个阿贝尔范畴。}

\VerdierSemanticAnchor{ch3d-2-1}
\subsection*{2.1. 初等态射。}

\VerdierSemanticAnchor{ch3d-2-1-1}
\noindent\textbf{2.1.1.}\quad
设 $X$ 是 $\mathcal A$ 的一个复形，$X^i$ 是它的 $i$ 次分量。设
$u\colon Y\longrightarrow X^i$ 是 $\mathcal A$ 的一个态射。构造如下交换图：

\VerdierSemanticAnchor{ch3d-2-1-1-diagram}
\begin{equation}
\begin{tikzcd}[column sep=1.45em,row sep=2.25em]
{} & {} & X'^{i-1} \arrow[r,"v_2"] \arrow[d,"u'"']
  & Y \arrow[d,"u"] \arrow[dr,"v_3"] & {} & {} \\
X^{i-3} \arrow[r,"{d_X^{i-3}}"]
  & X^{i-2} \arrow[r,"{d_X^{i-2}}"] \arrow[ur,"v_1"]
  & X^{i-1} \arrow[r,"{d_X^{i-1}}"]
  & X^i \arrow[r,"{d_X^i}"]
  & X^{i+1} \arrow[r,"{d_X^{i+1}}"]
  & X^{i+2}
\end{tikzcd}
\tag{2.1.1.1}
\end{equation}
其中方形

\VerdierSemanticAnchor{ch3d-2-1-1-cartesian-square}
\begin{center}
\begin{tikzcd}[column sep=4.2em,row sep=2.5em]
X'^{i-1} \arrow[r,"v_2"] \arrow[d,"u'"']
  & Y \arrow[d,"u"] \\
X^{i-1} \arrow[r,"{d_X^{i-1}}"'] & X^i
\end{tikzcd}
\end{center}
是笛卡尔方形，而态射 $v_1$ 与 $v_3$ 由下列等式定义：

\VerdierSemanticAnchor{ch3d-2-1-1-relations}
\[
v_2v_1=0,
\qquad
u'v_1=d_X^{i-2},
\qquad
v_3=d_X^iu.
\]
于是以 $X(u,i)$ 表示复形

\VerdierSemanticAnchor{ch3d-2-1-1-complex}
\begin{equation}
\begin{aligned}
X(u,i)={}&\cdots\longrightarrow X^{i-3}
\xrightarrow{d_X^{i-3}}X^{i-2}
\xrightarrow{v_1}X'^{i-1}
\xrightarrow{v_2}Y\\
&
\xrightarrow{v_3}X^{i+1}
\xrightarrow{d_X^{i+1}}X^{i+2}\longrightarrow\cdots,
\end{aligned}
\tag{2.1.1.2}
\end{equation}
并以 $p(u,i)\colon X(u,i)\longrightarrow X$ 表示分量如下的复形态射：

\VerdierSemanticAnchor{ch3d-2-1-1-projection}
\begin{equation}
\begin{aligned}
p(u,i)^j&=\operatorname{id},&&j\ne i-1,i,\\
p(u,i)^{i-1}&=u',\\
p(u,i)^i&=u.
\end{aligned}
\tag{2.1.1.3}
\end{equation}
\endgroup
